\section{Tính chất của kim loại và dãy hoạt động hóa học}

\subsection{Tính chất vật lí}

Kim loại thường có ánh kim, dẫn điện, dẫn nhiệt và có tính dẻo. Mức độ thể hiện các tính chất này khác nhau giữa các kim loại, vì vậy khi chọn vật liệu cần dựa vào yêu cầu sử dụng cụ thể.

\begin{ghinho}{Dãy hoạt động hóa học của kim loại}
\[
\ce{K, Na, Ca, Mg, Al, Zn, Fe, Pb, H, Cu, Ag, Au}
\]
Dãy hoạt động được dùng để dự đoán khả năng phản ứng của kim loại với nước, acid và dung dịch muối. Không nên chỉ học thuộc thứ tự; cần gắn thứ tự với hiện tượng và phương trình phản ứng.
\end{ghinho}

\subsection{Phản ứng với phi kim}

Nhiều kim loại phản ứng với oxygen, chlorine hoặc sulfur khi có điều kiện thích hợp.

\begin{tasks}(2)
\task $\ce{2Mg + O2 ->[t^\circ] 2MgO}$.
\task $\ce{3Fe + 2O2 ->[t^\circ] Fe3O4}$.
\task $\ce{2Fe + 3Cl2 ->[t^\circ] 2FeCl3}$.
\task $\ce{Fe + S ->[t^\circ] FeS}$.
\end{tasks}

\subsection{Phản ứng với nước và acid}

Một số kim loại hoạt động mạnh phản ứng với nước ở điều kiện thường. Các kim loại đứng trước H trong dãy hoạt động có thể phản ứng với một số acid loãng như \ce{HCl}, \ce{H2SO4} loãng để giải phóng hydrogen.

\[
\ce{2Na + 2H2O -> 2NaOH + H2 ^}
\]
\[
\ce{Ca + 2H2O -> Ca(OH)2 + H2 ^}
\]
\[
\ce{Zn + 2HCl -> ZnCl2 + H2 ^}
\]
\[
\ce{Fe + H2SO4 -> FeSO4 + H2 ^}
\]

\begin{notechem}
Khi dự đoán phản ứng với acid loãng, cần xét vị trí của kim loại so với H trong dãy hoạt động. Những trường hợp acid có tính oxi hóa mạnh nằm ngoài phạm vi quy tắc đơn giản này và chỉ xử lí khi đề bài cung cấp đủ dữ kiện.
\end{notechem}

\subsection{Phản ứng với dung dịch muối}

Một kim loại hoạt động mạnh hơn có thể đẩy kim loại yếu hơn ra khỏi dung dịch muối của nó, nếu phản ứng xảy ra trong môi trường nước phù hợp.

\[
\ce{Fe + CuSO4 -> FeSO4 + Cu}
\]
\[
\ce{Cu + 2AgNO3 -> Cu(NO3)2 + 2Ag}
\]

\begin{vidu}
Nhúng một thanh sắt sạch vào dung dịch \ce{CuSO4}. Nêu hiện tượng và viết phương trình hóa học.
\end{vidu}
\begin{loigiai}
Sắt hoạt động mạnh hơn đồng nên đẩy đồng ra khỏi dung dịch muối:
\[
\ce{Fe + CuSO4 -> FeSO4 + Cu}.
\]
Trên thanh sắt xuất hiện lớp đồng màu đỏ nâu; màu xanh của dung dịch \ce{CuSO4} nhạt dần.

\textbf{Vậy,} có đồng bám trên thanh sắt và dung dịch thay đổi màu do \ce{Cu^{2+}} bị tiêu thụ.
\end{loigiai}

\subsection{Suy ra thứ tự hoạt động từ thí nghiệm}

Khi đề không cho sẵn tên kim loại, hãy dùng các dấu hiệu:
\begin{enumerate}
\item kim loại có phản ứng với nước hay không;
\item kim loại có đẩy được H khỏi acid loãng hay không;
\item kim loại có đẩy được kim loại khác khỏi dung dịch muối hay không.
\end{enumerate}
Từ mỗi thí nghiệm, lập một quan hệ mạnh hơn -- yếu hơn rồi ghép các quan hệ lại.

\begin{tuluyen}
Ba kim loại A, B, C có các tính chất sau: A phản ứng với nước ở điều kiện thường; B không phản ứng với nước nhưng phản ứng với \ce{HCl}; C không phản ứng với \ce{HCl}. Sắp xếp mức độ hoạt động giảm dần của A, B, C.
\dapso{$A>B>C$.}
\begin{loigiaiBT}
A phản ứng với nước nên hoạt động mạnh. B đứng trước H vì phản ứng được với \ce{HCl}. C đứng sau H vì không phản ứng với \ce{HCl}.

\textbf{Vậy,} mức độ hoạt động giảm dần là $A>B>C$.
\end{loigiaiBT}
\end{tuluyen}
